% Sources: PAPER_MASTER_PLAN sec 9 (structure), RESULTS_027_CAMERA_READY.md
% (all numbers). Camera-ready figures: 500 maps, density 0.27, stage-2 base.

Teams of small autonomous robots are increasingly deployed for tasks in
which a single vehicle is insufficient: warehouse logistics, search and
rescue, environmental monitoring, and infrastructure inspection. In
cluttered environments each robot perceives only a fraction of its
surroundings --- on-board range sensors are limited in reach and are
routinely degraded by occlusion, reflective or absorptive surfaces, dust,
and transient hardware faults. \emph{Collaborative perception} addresses
this by letting teammates share what they sense, so that one robot's view
fills another's blind spot. In the navigation benchmark studied here --- ten
drones crossing a cluttered field while suffering sustained sensing
blackouts roughly one third of the time --- sharing is not a refinement but
a necessity: with collaborative perception 89.3\% of drones reach the goal,
without it only 45.9\% do, a gap of more than 43 percentage points under
identical conditions (Section~\ref{sec:results}).

The channel that carries this resilience, however, is itself an attack
surface. Shared observations must be fused into each robot's local obstacle
map, and safe fusion is conservative: if any teammate reports an obstacle
nearby, the robot should respect it. We show that this conservatism ---
implemented here by a per-direction \emph{minimum} over own and
received range estimates, the choice that respects any teammate's evidence
of a nearby obstacle --- hands a powerful lever to a compromised
insider. A single \emph{Byzantine} teammate that broadcasts fabricated
obstacles overrides even a fully-sighted receiver's own sensing, because a
fabricated near reading always wins the minimum. The attack is therefore
independent of the sensor-failure rate that motivates collaboration in the
first place: it succeeds against healthy sensors as easily as against
degraded ones. Two traitors among ten drones cut honest navigation success
by up to 22.5 percentage points, and a \emph{camouflage} variant that
extends a real obstacle into a free corridor is both more damaging and far
harder to detect than a blatant phantom wall.

The natural defense is cross-agent consistency checking: if a verifier is
positioned to see a reported obstacle and does not, the reporter loses
trust and is excluded from fusion. Under idealized sensing this neutralizes
the attack almost completely. It breaks down, however, in exactly the
conditions that make collaboration valuable. Under realistic ranging noise
two honest views of the same obstacle legitimately disagree by about
$\sqrt{2}\sigma$; once that exceeds the verifier's fixed tolerance, honest
neighbours are condemned as liars. Detection precision collapses from
$1.00$ to $0.23$, and merely enabling the filter destroys up to 27
percentage points of success in swarms containing no attacker at all:
\emph{the defense is worse than no defense}. Widening the tolerance with
$\sigma$ repairs the false accusations but opens the complementary hole ---
a camouflage attacker hides its fabrication \emph{inside} the widened band,
where no single-frame test can contradict it, and at the severest noise
level such a filter catches only 13\% of camouflage lies. The bind is
structural rather than a matter of tuning: on any single frame, a tolerance
tight enough to expose the lie is tight enough to condemn an honest
neighbour. That this difficulty is not hypothetical is borne out by an early
study of trust-modulated cooperative perception, which found that
cooperative fusion did not meaningfully improve accuracy over local
perception, with false detections inserted at high confidence among the
suspected causes~\cite{trupercept2020}.

Our central contribution closes this blind spot with time rather than with
looser thresholds. The per-frame \emph{offset vector} between a neighbour's
reported obstacle and the verifier's own view of that obstacle is zero-mean
for an honest neighbour --- two independent noise draws cancel in
expectation --- but retains a \emph{persistent bias} for a fabrication,
whose reported position is anchored to a lie rather than to the shared
ground truth. Aggregating this offset over a few tens of observations
separates the two populations cleanly. This use of time differs
fundamentally from recent temporal-consistency defenses for cooperative
perception, which compare a scene against its own recent past and so screen
for \emph{disruption}: once established, a persistent fabrication can
remain temporally self-consistent --- a signature such defenses do not
study --- whereas it leaves a persistent offset against the verifier's
\emph{own sensing}, which is exposed through accumulation over time
(Section~\ref{sec:related}). The resulting temporal trust rule
lifts camouflage detection recall from $0.13$ to $0.69$ and recovers up to
$+12.2$ percentage points of navigation success at the severest noise
level, while leaving attack-free swarms statistically untouched (the
measured cost is $-0.4$ percentage points with a confidence interval
spanning zero). Because each verdict is pairwise --- a robot tests each
neighbour against its own sensing and never counts votes --- the rule needs
no honest majority: its recovery remains statistically significant with up
to seven of ten robots lying, past the point where an honest robot's
neighbourhood is majority traitor, whereas the single-frame check is no
longer distinguishable from no defense there. Anticipating the strongest objection --- that a static
attacker is a strawman --- we also evaluate an adaptive attacker that knows
the defense and tunes its fabrication to evade it. The attacker faces a
\emph{stealth/harm bind}: as the phantom drifts far enough from real
geometry to block useful space, its offset signature grows and detection
recall rises with the harm; hiding from the filter forces the phantom onto
a real obstacle, where it blocks nothing new. We verify this bind at every
noise level and at one to three traitors.

This paper makes five contributions. \textbf{(i)}~We quantify the
resilience value of collaborative perception for learned multi-robot
navigation under sustained sensor dropout at a calibrated obstacle density,
with paired-bootstrap confidence intervals over 500-map suites.
\textbf{(ii)}~We characterize a minimum-fusion Byzantine vulnerability of
shared obstacle maps and show it is independent of the dropout rate.
\textbf{(iii)}~We show that naive consistency-based trust is not merely
weakened but \emph{destructive} under sensor noise. \textbf{(iv)}~We give a
noise-aware robust filter that restores safe operation and characterize its
single-frame camouflage limit. \textbf{(v)}~We introduce a temporal
offset-bias trust rule that recovers this limit at no measured cost to
honest swarms and, being pairwise, retains a significant recovery with up to
seven of ten robots lying --- past the loss of an honest local majority ---
and we validate it against an adaptive, defense-aware attacker
across one to three traitors, demonstrating a stealth/harm bind that makes
the recovery difficult to game.

The remainder of the paper is organized as follows.
Section~\ref{sec:related} positions the work against collaborative
perception security, temporal spoof detection, and Byzantine multi-robot
systems. Section~\ref{sec:methods} details the environment, the learned
navigation policy, the threat model, and the three defenses.
Section~\ref{sec:setup} defines metrics and the statistical protocol.
Section~\ref{sec:results} reports results, Section~\ref{sec:discussion}
discusses limitations honestly, and Section~\ref{sec:conclusion} concludes.
